\documentclass[11pt]{article}

\usepackage[margin=1in]{geometry}
\usepackage{amsmath, amssymb}
\usepackage{algorithm}
\usepackage{algpseudocode}
\usepackage{booktabs}
\usepackage{hyperref}
\usepackage{braket}

\title{CSE 4412: Introduction to Quantum Computing - Week 2 Notes}
\author{}
\date{}

\begin{document}
\maketitle
\section*{Postulate 2: measurements}
This is how we can extract information from a quantum system. We can measure a quantum state $\ket{\psi}$ by applying a measurement operator $M$ to it. The measurement operator is a set of projectors $\{P_i\}$ that correspond to the possible outcomes of the measurement. When we apply the measurement operator to the state, we get a new state that is one of the possible outcomes, and the probability of getting that outcome is given by the Born rule: $P(i) = \bra{\psi} P_i \ket{\psi}$. 
\vspace{1cm}

where $\ket{+} = \frac{1}{\sqrt{2}}(\ket{0} + \ket{1})$ and $\ket{-} = \frac{1}{\sqrt{2}}(\ket{0} - \ket{1})$. where $\mathbb{P}_i$ is the probability of measuring the state $\ket{\psi}$ in the state $\ket{i}$, and $\bra{\psi} P_i \ket{\psi}$ is the inner product of the state $\ket{\psi}$ with the projector $P_i$ applied to it. Given that thye are bases we can rewrite $\ket{0}$ and $\ket{1}$ in terms of $\ket{+}$ and $\ket{-}$ as follows: 
\[
\ket{0} = \frac{1}{\sqrt{2}}(\ket{+} + \ket{-}) \\ {\text{, and }} \\
\ket{1} = \frac{1}{\sqrt{2}}(\ket{+} - \ket{-})
\]

\section*{Postulate 4: How to model two or more states}
We can use tensor product to model two or more quantum states. If we have two quantum states $\ket{\psi_1}$ and $\ket{\psi_2}$, we can combine them into a single state using the tensor product:

\[
\ket{\psi_1} \otimes \ket{\psi_2}
\]

\end{document}
